\paragraph{Forced outside edges and an optimal selector.}
Let $G_s,G_t$ be full fibre components with arbitrary Hamilton cycles
$C,D$ and prescribed pools $P,Q$. Put
$U=V(G_s)\cap V(G_t)$ and $O=V(G_s)\setminus U$.
The graphs have the same recolouring edges on common endpoints.
Let $F=(P\cup Q)\cap E(G_s[U])$ and $f(v)=d_F(v)$.
Let $J$ consist of those edges of $G_s[O]$ having an endpoint of
degree two in the full source component. Every source Hamilton cycle
contains $J$, so $d_J(w)\le2$ and
$c(w)=2-d_J(w)\ge0$ for $w\in O$.

For $S\subseteq U$, define
\[
 M_J(S)=\sum_{v\in S}(4-d_t(v)-f(v))
       -\sum_{w\in O}\min\{|N_s(w)\cap S|,c(w)\}.
\]
Then $(C\setminus P)\cap(D\setminus Q)$ contains at least
$\max\{0,\lceil M_J(S)/2\rceil\}$ edges.
Indeed, the cycle incidences from $S$ to an outside vertex $w$ are
at most $|N_s(w)\cap S|$ and at most $c(w)$, because its two source
cycle incidences already include $d_J(w)$ edges within $O$.
The lost common incidences due to the pools are at most
$\sum_{v\in S}f(v)$. At a common vertex $v$, if $\ell(v)$ source
cycle incidences are lost, $2-\ell(v)$ safe source edges remain among
the $d_t(v)$ target edges. The target cycle uses two of these target
edges, so there are at least $4-d_t(v)-\ell(v)$ common safe
incidences at $v$. Summing and counting each common edge at most twice
gives the bound. This proof uses neither selected cycle edge set.

Suppose the boundaries differ in one non-star coordinate $j$, from
$a$ to $b$. A source vertex belongs to $O$ precisely when its colour
at $j$ is $b$. Thus each $v\in U$ has at most one source neighbour
in $O$: the word obtained by changing its $j$-th colour to $b$.
Conversely each $w\in O$ has at most two source neighbours in $U$,
since their $j$-th colours avoid both $a$ and $b$.

Put $r(v)=4-d_t(v)-f(v)$. The set
\[
 S_* = \{v\in U:r(v)>0\}
\]
maximizes $M_J(S)$. Removing a vertex with $r(v)\le0$ never
decreases the objective. Adding a vertex with $r(v)>0$ adds an
integer reward of at least one and raises the outside-capacity cost
by at most one, because it has at most one outside neighbour.
Therefore all such vertices may be added. Since a target Hamilton
cycle exists, $d_t(v)\ge2$; positive reward implies $d_t(v)\le3$.
This selector proof requires a single non-star boundary edit;
the conditional bound for a specified $S$ does not.

If $J$ is empty, the outside capacity is two, and
$|N_s(w)\cap S|\le2$ makes every capacity truncation inactive.
Writing $d_c(v)=|N_s(v)\cap N_t(v)|$, the optimized margin is
\[
 M_0(S_*)=\sum_{v\in U}
   \max\{0,4-d_s(v)-d_t(v)+d_c(v)-f(v)\}.
\]
Indeed $d_s(v)-d_c(v)=|N_s(v)\cap O|\le1$, so the selected
vertices have nonnegative summands; omitted vertices have
nonpositive summands. Consequently the optimized original aggregate
and the positive node bound have exactly the same coverage.
The forced-edge capacities can strengthen this bound without
requiring any individual summand to be positive.

On the fixed 230 ordinary phase6 interfaces, the original all-low-degree
selector is positive in 97 cases. The optimal selector is positive in
125 cases; the forced-edge capacities remain positive in the same
125 cases. They raise the guaranteed edge count from one to two for
$00123$ with $(j,q)=(2,3)$ and $(4,1)$, both normalizing to $00121$.
The common outside vertex $31301$ has two forced internal outside
edges, hence zero remaining capacity, saving one lost incidence in
each case. The largest guaranteed count is six. These are finite
input checks, not universal positive-coverage claims; hard port-seed
components are not substituted for whole-component Hamilton cycles.
